\input{../tex-inputs/latex-leseansicht-vorspann}

\section[Gustav und Max Schwarzkopf an Arthur Schnitzler, 15. 9. 1909]{L04272 Gustav und Max Schwarzkopf an Arthur Schnitzler, 15. 9. 1909}
\nopagebreak\mylabel{L04272v}
\rehead{ }\normalsize\beginnumbering\briefempfaengerindex{Schnitzler, Arthur@\textsc{Schnitzler, Arthur}!zzzSchwarzkopf, Max@\emph{von Max Schwarzkopf}!1909-09-153@{15. 9. 1909}|(be}\briefempfaengerindex{Schnitzler, Arthur@\textsc{Schnitzler, Arthur}!zzzSchwarzkopf, Gustav@\emph{von Gustav Schwarzkopf}!1909-09-153@{15. 9. 1909}|(be}
\toendnotes[C]{\smallbreak\pagebreak[2]}
\correspDesc{Versand  durch Gustav Schwarzkopf, Max Schwarzkopf am 15. 9. 1909 in Wien
\newline{}Erhalt  durch Arthur Schnitzler im Zeitraum [15. 9. 1909 – 18. 9. 1909?] in Wien}\toendnotes[C]{\smallbreak}
\Standort{CUL, Schnitzler, B 96a.}
\physDesc{Briefkarte, 402 Zeichen
\newline{}Handschrift Gustav Schwarzkopf: schwarze Tinte, deutsche Kurrent
\newline{}Handschrift Max Schwarzkopf: schwarze Tinte, deutsche Kurrent
\newline{}Schnitzler: mit Bleistift Vermerk: »\textsc{Schwarzkopf}« }\toendnotes[C]{\smallbreak}
\pstart
           \raggedleft{}{\pb}15. IX 09\pend
           \vspace{0.5em}
\pstart
           Lieber Arthur, ich freue mich, daß Ihr Wunſch erfüllt wurde, daß
               alles gut gegangen iſt und gratulire. Ihrer Frau\pwindex{Schnitzler, Olga 17.\,1.\,1882 Wien – 13.\,1.\,1970 Lugano@\textsc{Schnitzler, Olga} (17.\,1.\,1882 Wien – 13.\,1.\,1970 Lugano), \emph{Schauspielerin, Sängerin}|pwv} und Ihnen herzlichſt. Pünktlich ist ſie\pwindex{Cappellini, Lili 13.\,9.\,1909 Wien – 26.\,7.\,1928 Venedig@\textsc{Cappellini, Lili} (13.\,9.\,1909 Wien – 26.\,7.\,1928 Venedig)|pwv} übrigens geweſen, die junge Dame\pwindex{Cappellini, Lili 13.\,9.\,1909 Wien – 26.\,7.\,1928 Venedig@\textsc{Cappellini, Lili} (13.\,9.\,1909 Wien – 26.\,7.\,1928 Venedig)|pwv}, das muß man ihr laſſen.
                  {\pb}We{\geminationn}
               mein Halsweh es geſtattet, werde ich mir erlauben mich dem Fräulein Lili\pwindex{Cappellini, Lili 13.\,9.\,1909 Wien – 26.\,7.\,1928 Venedig@\textsc{Cappellini, Lili} (13.\,9.\,1909 Wien – 26.\,7.\,1928 Venedig)|pw} vorzuſtellen.\pend
           
\pstart
           Ihre Frau\pwindex{Schnitzler, Olga 17.\,1.\,1882 Wien – 13.\,1.\,1970 Lugano@\textsc{Schnitzler, Olga} (17.\,1.\,1882 Wien – 13.\,1.\,1970 Lugano), \emph{Schauspielerin, Sängerin}|pwv} und Sie herzlichſt{\\[\baselineskip]}grüßend{\\[\baselineskip]}Ihr{\\[\baselineskip]}\spacefill\mbox{Gustav.}\pend
           \leftskip=0em{}\selectlanguage{ngerman}\vspace{1em}{\vspace{1\baselineskip}}
\pstart
           {[}hs. Schwarzkopf:{]} Auch von mir die herzlichſten Glückwünſche.\pend
           
\pstart
           Ihr{\\[\baselineskip]}\spacefill\mbox{MaxSchwarzk}\pend
           \leftskip=0em{}\selectlanguage{ngerman}\endnumbering\briefempfaengerindex{Schnitzler, Arthur@\textsc{Schnitzler, Arthur}!zzzSchwarzkopf, Max@\emph{von Max Schwarzkopf}!1909-09-153@{15. 9. 1909}|)be}\briefempfaengerindex{Schnitzler, Arthur@\textsc{Schnitzler, Arthur}!zzzSchwarzkopf, Gustav@\emph{von Gustav Schwarzkopf}!1909-09-153@{15. 9. 1909}|)be}\mylabel{L04272h}
\begin{anhang}
\end{anhang}\newcommand{\dateiname}{L04272}\newcommand{\titel}{Gustav und Max Schwarzkopf an Arthur Schnitzler, 15. 9. 1909}\newcommand{\editorInnen}{Herausgegeben von Jahnke, SelmaMüller, Martin Anton}\input{../tex-inputs/latex-leseansicht-abspann}
